\documentclass[11pt]{article}
\usepackage[T1]{fontenc}
\usepackage[utf8]{inputenc}
\usepackage{lmodern}
\usepackage[margin=1in]{geometry}
\usepackage{amsmath,amssymb,amsthm,mathtools}
\usepackage{booktabs,array,longtable}
\usepackage{enumitem}
\usepackage{hyperref}
\usepackage{xurl}

\IfFileExists{../paper/release_info.tex}{
  \newcommand{\OPHPaperReleaseID}{r1520}
\newcommand{\OPHPaperReleaseDate}{July 9, 2026}
\newcommand{\OPHPaperReleaseBanner}{%
\begin{center}
\small\textbf{Paper release:} \texttt{\OPHPaperReleaseID}\quad\textbf{Released:} \OPHPaperReleaseDate
\end{center}
}

}{
  \newcommand{\OPHPaperReleaseDate}{\today}
  \newcommand{\OPHPaperReleaseBanner}{Draft}
  \newcommand{\OPHPaperReleaseID}{local-cosmology-staging}
}

\hypersetup{colorlinks=true,linkcolor=blue,citecolor=blue,urlcolor=blue}
\setlength{\emergencystretch}{3em}
\setlength{\tabcolsep}{4pt}
\renewcommand{\arraystretch}{1.08}

\newtheorem{theorem}{Theorem}[section]
\newtheorem{definition}[theorem]{Definition}
\newtheorem{assumption}[theorem]{Assumption}
\newtheorem{target}[theorem]{Target}
\theoremstyle{remark}
\newtheorem{remark}[theorem]{Remark}
\theoremstyle{plain}

\newcommand{\pathref}[1]{\nolinkurl{#1}}
\newcommand{\SimSourceRepo}{\url{https://github.com/muellerberndt/oph-physics-sim}}

\title{OPH Cosmological Vacuum and Structure Formation:\\
Finite Transfer Gates, Fluctuation Ensembles, and Proto-Object Receipts}
\author{B.\ M\"uller}
\date{\OPHPaperReleaseDate}

\begin{document}
\maketitle
\OPHPaperReleaseBanner

\begin{abstract}
OPH cosmology needs a sharp distinction between a field that looks noisy and a
field that deserves to be called a physical vacuum. This paper sets that
boundary. A settled simulator state, a repair relaxation, or a plausible visual
fluctuation is only a diagnostic unless it comes from OPH source data with a
proper transfer construction and refinement check.

The same boundary organizes structure formation. Vacuum-like fields, primordial
seeds, and observable sky predictions are different claims. The paper keeps
those layers separate so a useful simulation cannot be mistaken for a completed
cosmological prediction.
\end{abstract}

\section{Scope}

This is a technical companion outside the release bundle. It coordinates the
vacuum and fluctuation boundary for finite-source cosmology, screen
microphysics, the compact theorem stack, and OPH-FPE simulator outputs.

\section{Shared Finite Quotient Theorem Surface}

\paragraph{Finite quotient ensemble theorem surface.}
Fix a finite regulator \(r\).  The physical presentation space is \(\Sigma_r\), the
presentation redundancy groupoid is \(\Gamma_r\), and the finite physical quotient is
\[
Q_r=\Sigma_r/\Gamma_r,\qquad \pi_r:\Sigma_r\to Q_r.
\]
The quotient removes nonphysical presentation data: gauge representatives, port relabelings, mesh
labels, shard or worker identifiers, queue order, repair schedule identifiers, retry counters,
timestamps unless declared semantic, hidden carrier coordinates, and inert ancillary labels.  If
only settled configurations carry probability, the probability space is the normal-form subset
\[
N_r=n_r(Q_r).
\]
The map \(n_r\) is a normal-form map, not a probability law. Any promoted physical branch inherits
this firewall: it must declare the quotient-intrinsic source law or action before a normal form can
be read as a selection or prediction claim.

\paragraph{Observable algebras and reference states.}
In the finite classical case the quotient observable algebra is
\[
\mathcal O_r=\ell^\infty(Q_r).
\]
In the finite quantum case the physical algebra is a declared quotient algebra
\(\mathcal A^{\rm phys}_r\) with state
\[
\omega_r(A)=\operatorname{Tr}(\rho_r A).
\]
When the reference object is obtained from a finite lifted carrier, the load-bearing data are not
an abstract groupoid cardinality alone.  They are a tracially pointed quotient
\[
\left(\mathcal A^{\rm phys}_{r,b},\tau^0_{r,b}\right),
\qquad
\mathcal A^{\rm phys}_{r,b}
=
z_{r,b}B(\widetilde{\mathcal H}_r)^{G_r}z_{r,b},
\]
where \(U_r:G_r\to U(\widetilde{\mathcal H}_r)\) is the compact gauge action and \(z_{r,b}\) is
the central projection for the declared boundary or superselection sector.  The reference trace is
\[
\tau^0_{r,b}(A)
=
\frac{\operatorname{Tr}_{\widetilde{\mathcal H}_r}(A)}
{\operatorname{Tr}_{\widetilde{\mathcal H}_r}(z_{r,b})}.
\]
If
\[
z_{r,b}\widetilde{\mathcal H}_r
\cong
\bigoplus_\alpha V_\alpha\otimes M_\alpha,
\]
with \(d_\alpha=\dim V_\alpha\) and \(m_\alpha=\dim M_\alpha\), then
\[
\mathcal A^{\rm phys}_{r,b}
\cong
\bigoplus_\alpha I_{V_\alpha}\otimes B(M_\alpha),
\qquad
p_{r,\alpha}
=
\frac{d_\alpha m_\alpha}{\sum_\beta d_\beta m_\beta}.
\]
These are the induced central-sector weights only after the carrier representation and boundary
sector have been fixed.

\paragraph{OPH quotient ensemble.}
An OPH quotient ensemble is specified by a quotient-intrinsic base weight and action
\[
m_r:Q_r\to \mathbb R_{>0},
\qquad
S_r:Q_r\to \mathbb R\cup\{+\infty\},
\]
and
\[
w_r(q)=m_r(q)e^{-S_r(q)},\qquad
Z_r=\sum_{q\in Q_r}w_r(q),\qquad
\mu_r(q)=Z_r^{-1}w_r(q).
\]
Equivalently, one may state an intrinsic projective prior \(\nu_r\) on \(Q_r\) and set
\(\mu_r=(n_r)_\#\nu_r\).  Uniform quotient counting, uniform representative counting pushed to the
quotient, groupoid weights, and tracial central-sector weights are different physical claims.  The
paper must declare which one is being used.

\paragraph{Normal-form projector non-selection.}
For any retraction \(N:Q\to Q_{\rm nf}\) onto a subset \(Q_{\rm nf}\subseteq Q\), that is, any map
whose restriction to \(Q_{\rm nf}\) is the identity, the induced map on laws
\[
\mathcal C_Q(\mu)=N_\#\mu
\]
is idempotent:
\[
\mathcal C_Q^2=\mathcal C_Q.
\]
Every law supported on \(Q_{\rm nf}\) is fixed; both statements use the retraction property.  Therefore settlement or canonicalization never
selects a unique physical probability law by itself.

\paragraph{Selection-gap corollary.}
Let \(X\subseteq Q_{\rm nf}\) be a finite set of quotient-normal candidates distinguished by
visible invariants.  Normal-form data determine \(X\) and its quotient-visible invariants, but they
do not choose a member of \(X\).  If two laws \(\mu,\nu\) are supported on \(X\) and concentrate on
different candidates, both are fixed by \(\mathcal C_Q\).  Unique sector selection therefore
requires source data: an intrinsic action with a unique minimizer, a declared physical ensemble, or
a refinement-stable gap certificate.  A defect or holonomy classification can classify possible
sectors without choosing the physical sector, and a contraction or repair generator can certify
convergence toward a declared target without creating the target law.

\paragraph{Finite MaxEnt quotient ensemble.}
For finite \(Q\), positive \(m\), and quotient observables \(F_1,\ldots,F_k\), maximizing
\[
\mathcal H_m(\nu)=-\sum_q\nu(q)\log\frac{\nu(q)}{m(q)}
\]
subject to
\[
\sum_q\nu(q)=1,\qquad \sum_q\nu(q)F_a(q)=c_a
\]
has the full-support solution, when the feasible full-support surface is nonempty,
\[
\mu(q)=
\frac{m(q)\exp[-\sum_a\theta_aF_a(q)]}{Z(\theta)}.
\]
Boundary optima obey the same formula after restricting to their support.  On a finite
noncommutative quotient algebra with faithful reference state \(\sigma_r\),
\[
\rho_r=
\frac{\exp(\log\sigma_r-\sum_a\theta_aF_{r,a})}
{\operatorname{Tr}\exp(\log\sigma_r-\sum_a\theta_aF_{r,a})}.
\]
The finite constraint ledger must name every \(F_{r,a}\), its units and support, the target
expectation and source, sector or zero-mode treatment, refinement transformation, and proof that no
run output or observational output entered the source definition.

\paragraph{Refinement compatibility and RG closure.}
For \(s\succeq r\), let \(c_{sr}:Q_s\to Q_r\) be the physical coarse map.  Exact compatibility of
weighted ensembles is equivalent to the fiber-sum identity
\[
\sum_{q':\,c_{sr}(q')=q}m_s(q')e^{-S_s(q')}
=
\alpha_{sr}m_r(q)e^{-S_r(q)}
\]
for a constant \(\alpha_{sr}>0\) independent of \(q\).  Then
\[
(c_{sr})_\#\mu_s=\mu_r.
\]
If the one-step defects are
\[
\delta_{k+1,k}
=
\left\|(c_{k+1,k})_\#\mu_{k+1}-\mu_k\right\|_{\mathrm{TV}},
\]
then
\[
\left\|(c_{nr})_\#\mu_n-\mu_r\right\|_{\mathrm{TV}}
\le
\sum_{k=r}^{n-1}\delta_{k+1,k}.
\]
For exponential-family refinement, exact closure requires the fine conditional free energy
\[
G_{sr,\theta}(q_r)
=
-\log\mathbb E_{m_s^0}\left[
\exp[-\theta\cdot F_s(Q_s)]\mid c_{sr}(Q_s)=q_r
\right]
\]
to equal \(\kappa_{sr}(\theta)+R_{sr}(\theta)\cdot F_r(q_r)\).  If the residual is uniformly
bounded by \(\varepsilon\), the induced total-variation defect is bounded by \(\tanh\varepsilon\).

\paragraph{Implementation invariance and representative lifting.}
If implementations \(A,B\) have quotient bijections \(h_r:Q_r^A\to Q_r^B\) satisfying
\[
m_r^B(h_rq)=m_r^A(q),\qquad
S_r^B(h_rq)=S_r^A(q),
\qquad
h_r\circ c_{sr}^A=c_{sr}^B\circ h_s,
\]
then
\[
(h_r)_\#\mu_r^A=\mu_r^B.
\]
For tracially pointed quantum quotients the corresponding equivalence is a trace-preserving
quotient equivalence.  It is invariant under unitary intertwiners preserving the gauge action and
sector, and under inert trivial ancillas \(A\mapsto A\otimes I_{\rm anc}\).  It is not invariant
under arbitrary changes of gauge-representation multiplicities.

If an implementation stores representatives, a representative-level law must be a conditional lift
\[
\widetilde\mu_r(x)=\mu_r(\pi_r x)\kappa_r(x\mid \pi_r x),
\qquad
\sum_{x:\pi_r(x)=q}\kappa_r(x\mid q)=1.
\]
Then \((\pi_r)_\#\widetilde\mu_r=\mu_r\).  Uniform representative sampling yields orbit-size
weights and is physical only if representative counting is the declared base measure.

\paragraph{Quotient-lumpable kernels and sampler correctness.}
A representative kernel \(\widetilde P(x,y)\) descends to \(Q_r\) only when
\[
P_Q(q,q')
=
\sum_{y:\pi(y)=q'}\widetilde P(x,y)
\]
is independent of the chosen representative \(x\in\pi^{-1}(q)\).  For
\(w(q)=m(q)e^{-S(q)}\) and proposal \(R(q,q')\) with reciprocal support, the
Metropolis--Hastings acceptance rule
\[
a(q,q')=\min\left\{1,
\frac{w(q')R(q',q)}{w(q)R(q,q')}
\right\}
\]
gives detailed balance
\[
\mu(q)R(q,q')a(q,q')=\mu(q')R(q',q)a(q',q).
\]
Repair-informed proposals must include the Hastings asymmetry term; otherwise the stationary law is
generically changed.

\paragraph{Repair generators are not selectors.}
A repair generator of the form
\[
L_{\rm rep}=\sum_C c_C(I-E_C)
\]
is a relaxation or sampling object after a law has been selected.  Conditional expectations
\(E_C\) are defined on \(L^2(X_r,\pi_r)\), so the reference law \(\pi_r\) is input.  On
overlapping collars the expectations need not commute.  The correct finite gap certificate is the
Poincare constant
\[
\kappa_r
=
\inf_{f\perp 1}
\frac{\sum_C\|(I-E_C)f\|^2}{\|f\|^2}.
\]
If local fiber rates have a positive lower bound \(\gamma_*\), then
\[
L_{\rm rep}\ge \gamma_*\kappa_r(I-P_0).
\]
Finite repair completeness gives \(\kappa_r>0\) at fixed regulator.  A uniform refinement lower
bound \(\inf_r\kappa_r>0\) is a separate theorem or receipt.

\paragraph{Finite evidence accuracy.}
For bounded coarse observables \(O\), if
\[
\|\widehat\mu_s-\mu_s\|_{\mathrm{TV}}\le\epsilon_{\rm samp}
\]
and the refinement defects sum to \(\epsilon_{\rm ref}\), then
\[
\left|\mathbb E_{\widehat\mu_s}[O\circ c_{sr}]-\mathbb E_{\mu_r}[O]\right|
\le
2\|O\|_\infty(\epsilon_{\rm samp}+\epsilon_{\rm ref}).
\]
Continuum-facing observables require a realization map and correlation Cauchy bound in addition to a
finite histogram.

\paragraph{Vacuum promotion gate.}
A stationary sampler is not a physical vacuum.  For any faithful target law one can build a
positive transfer operator with that law as ground state, so positivity alone is not a selector.
Vacuum promotion requires source Euclidean slab data
\[
\mathfrak S_r^E=(Q_r,m_r^0,J_r,V_r,a_{t,r})
\]
whose conductance \(J_r(q,q')=J_r(q',q)\ge0\), local potential \(V_r\), and slab thickness
\(a_{t,r}\) are derived without using the target law or sampler output.  With connected event graph,
\[
(H_r^Ef)(q)
=
\frac{1}{m_r^0(q)}
\sum_{q'}J_r(q,q')\bigl(f(q)-f(q')\bigr)
+
V_r(q)f(q)
\]
is self-adjoint and bounded below on \(L^2(Q_r,m_r^0)\); its finite Feynman--Kac semigroup is
positivity improving.  Perron--Frobenius gives a unique positive normalized ground state
\(\Omega_r\), and the finite vacuum law is
\[
\mu_r^{\rm vac}(q)=|\Omega_r(q)|^2m_r^0(q).
\]
For \(T_r=e^{-a_{t,r}(H_r^E-E_{0,r})}\), the Doob kernel is stochastic and detailed-balanced with
\(\mu_r^{\rm vac}\).  Continuum promotion additionally requires reflection positivity or
equivalent reconstruction plus refinement compatibility of the transfer family.

\paragraph{Primordial and cosmological prediction firewall.}
A screen covariance is not automatically a primordial curvature spectrum.  The map
\[
C_\ell
=
4\pi
\int_0^\infty
\frac{dk}{k}
\Delta_\zeta^2(k)j_\ell^2(k\chi_\star)
\]
has a null space unless a radial prior, finite parametric family, or source-stress bridge is
declared.  OPH primordial promotion requires the source-only stress, single-clock, entropy-repair,
curvature-evolution, adiabatic-mode, phase-coherence, screen-to-radial-lift, radial-null-space, and
forward-projection receipts. Observable CMB comparison also requires declared source, solver,
dataset, covariance, nuisance, data-use, and pooled-reducer provenance.

\paragraph{Claim tiers and required receipts.}
Every ensemble-facing run records its ensemble id, claim tier, regulator, representative
schema, gauge action, canonicalizer, base measure, action coefficients, coarse maps,
zero-mode projector, amplitude convention, sampler, smoothing policy, source provenance, and
explicit nonclaims. The seed belongs to the run receipt rather than the ensemble definition. The
claim tiers are
\[
\begin{array}{ll}
E0:&\text{seed noise, proposal noise, repair jitter},\\
E1:&\text{conventional reference ensemble},\\
E2:&\text{OPH-native quotient ensemble},\\
E3:&\text{OPH vacuum},\\
E4:&\text{OPH primordial field},\\
E5:&\text{observable cosmological prediction}.
\end{array}
\]
The evidence bundle must keep separate receipts for stationary-law schedule invariance, detailed balance
of the aggregate kernel, and pathwise partition invariance.  Deterministic replay of semantic
random streams or a canonical serial chain is useful, but it is not pathwise partition invariance.
Smoothing must preserve raw coefficients, raw spectra, smoothing kernels, smoothed coefficients,
smoothed spectra, and hashes of each stage; it is not part of \(S_r\) unless explicitly declared.


\section{Vacuum Claim Boundary}

The strongest allowed simulator labels are:
\[
E0:\hbox{ seed noise or repair jitter},\qquad
E1:\hbox{ conventional reference ensemble}.
\]
An OPH-native vacuum label requires an E3 receipt.  The source data must include
\[
\mathfrak S_r^E=(Q_r,m_r^0,J_r,V_r,a_{t,r})
\]
with \(J_r\), \(V_r\), and \(a_{t,r}\) derived from OPH source structure rather than from a target
law or a desired visual spectrum.  Accepted repair moves are generally directed and cannot simply
be reused as reversible Euclidean conductances.

\section{Structure-Formation Targets}

\begin{longtable}{@{}p{0.28\textwidth}p{0.62\textwidth}@{}}
\toprule
Object & Required source-side label\\
\midrule
\endhead
\bottomrule
\endfoot
Fluctuating screen field & E0 or E1 unless an E2 quotient ensemble receipt passes.\\
OPH vacuum field & E3 only after source Euclidean slab, ground-state, Doob, reflection-positivity,
and transfer-refinement receipts.\\
Proto-object worldlines & Diagnostic until production matter/particle and observer-facing object
population receipts pass.\\
Growth seeds & E4 only after the primordial bridge, radial lift, null-space, positivity, and
forward-projection receipts pass.\\
LSS comparison & E5 only after frozen likelihood/data contracts pass in the data-contract paper.\\
\end{longtable}

\section{Simulator Contract}

The simulator should emit the following receipt names even when they fail:
\begin{itemize}[leftmargin=2em]
\item \pathref{TRACIALLY_POINTED_QUOTIENT_RECEIPT}
\item \pathref{SOURCE_EUCLIDEAN_SLAB_RECEIPT}
\item \pathref{TRANSFER_REFINEMENT_RECEIPT}
\item \pathref{STATIONARY_LAW_SCHEDULE_INVARIANCE_RECEIPT}
\item \pathref{PATHWISE_PARTITION_INVARIANCE_RECEIPT}
\end{itemize}
False receipts are part of the scientific result.  They prevent a reference baseline or rendering
field from being mislabeled as an OPH vacuum.

\section{Unclosed Claim Boundaries}

\begin{enumerate}[leftmargin=2em]
\item Derive the reversible source conductances \(J_r\), potential \(V_r\), and slab thickness
\(a_{t,r}\) from finite OPH source data.
\item Prove reflection positivity or an equivalent reconstruction theorem on the selected branch.
\item Relate passed E3 vacuum receipts to E4 primordial growth seeds without using CMB or LSS data.
\item Specify proto-object/worldline receipts that remain distinct from production particle claims.
\end{enumerate}

\end{document}
